\documentclass[11pt,a4paper]{article}
\usepackage{fontspec}
\setmainfont{DejaVu Sans}[Scale=0.90]
\usepackage{amsmath,amssymb}
\usepackage[margin=2.2cm]{geometry}
\usepackage{booktabs,array,tabularx}
\usepackage{enumitem}
\usepackage{titlesec}
\usepackage{parskip}
\titlelabel{\thetitle.\quad}
\titleformat{\section}{\normalfont\large\bfseries}{\thesection.}{0.6em}{}
\titleformat{\subsection}{\normalfont\normalsize\bfseries}{\thesubsection}{0.6em}{}
\titlespacing*{\section}{0pt}{15pt}{5pt}
\titlespacing*{\subsection}{0pt}{11pt}{4pt}
\setlist[itemize]{leftmargin=1.4em,itemsep=1pt,topsep=3pt}
% Preserve fonts/margins; permit a final line-breaking pass for long prose.
\setlength{\emergencystretch}{2em}
\begin{document}
\begin{center}
{\Large\bfseries De Boltzmann à Guyer--Krumhansl}\\[4pt]
{\normalsize Collisions, approximation de Callaway, hiérarchie de moments, fermeture}\\[2pt]
{\small Notes de partie I --- 12 septembre 2026}
\end{center}

\vspace{4pt}


Ces notes établissent la chaîne conduisant de l'équation de transport des phonons à une équation
macroscopique du flux de chaleur. Les objets employés --- fonction de distribution, moments,
quasi-moment, distribution déplacée --- sont ceux des notes précédentes.

\section{L'équation de transport}

\subsection{Forme générale}

Le nombre de phonons d'un mode donné, au voisinage d'un point, varie pour deux raisons : les
phonons se déplacent, et ils entrent en collision. Le bilan s'écrit
\[ \boxed{\;\frac{\partial f}{\partial t} + \mathbf{v}_g \cdot \nabla_\mathbf{r} f
= \left(\frac{\partial f}{\partial t}\right)_{\!\mathrm{coll}}\;} \]

Le terme de gauche décrit le transport libre : un phonon de vitesse de groupe $\mathbf{v}_g$ quitte
une région et en rejoint une autre. Le terme de droite décrit les collisions.

Les phonons ne portant pas de charge, aucune force extérieure n'agit sur eux dans un cristal
uniforme ; le terme en $\nabla_\mathbf{k} f$ qui figure dans l'équation de Boltzmann d'un gaz
classique est donc absent.

\subsection{Pourquoi cette équation n'est pas utilisable telle quelle}

C'est le point de départ cinétique sous les hypothèses semi-classiques retenues. Sa résolution directe reste coûteuse.

L'inconnue $f(\mathbf{r}, \mathbf{k}, t)$ dépend de sept variables --- trois de position, trois de
vecteur d'onde, une de temps --- et porte de surcroît un indice de branche.

Le terme de collision est une intégrale sur tous les couples de modes susceptibles d'interagir, et
il est quadratique en $f$ pour les processus à trois phonons.

Les coefficients de couplage entre modes exigent la connaissance des constantes de force
anharmoniques du cristal.

Toute la suite consiste donc à remplacer cette équation par un petit nombre d'équations portant sur
des moments, au prix d'hypothèses explicites.

\section{Les collisions}

\subsection{Processus normaux}

Un processus à trois phonons combine deux phonons en un troisième, ou l'inverse. La conservation de
l'énergie impose $\omega_1 + \omega_2 = \omega_3$. Pour le quasi-moment, deux cas se présentent.

Lorsque
\[ \mathbf{k}_1 + \mathbf{k}_2 = \mathbf{k}_3 , \]
le quasi-moment total est conservé. Le processus est dit \emph{normal}.

Un tel processus redistribue l'énergie entre modes, mais laisse inchangée la somme
$\sum \hbar\mathbf{k}$. Comme le flux de chaleur est proportionnel à cette somme pour une dispersion
linéaire, \textbf{un processus normal ne dégrade pas le flux}. Il ne crée donc aucune résistance
thermique.

Son rôle n'est pas nul pour autant : il établit l'équilibre interne du gaz de phonons, et pousse la
distribution vers la forme déplacée.

\subsection{Processus umklapp}

Le quasi-moment n'étant défini qu'à un vecteur du réseau réciproque près, la conservation peut
s'écrire
\[ \mathbf{k}_1 + \mathbf{k}_2 = \mathbf{k}_3 + \mathbf{G} , \qquad \mathbf{G} \neq \mathbf{0} , \]
où $\mathbf{G}$ est un vecteur du réseau réciproque. Le phonon résultant est ramené dans la première
zone de Brillouin, et sa direction peut être opposée à celle des phonons incidents.

Ces processus portent le nom allemand d'\emph{umklapp}, retournement. Ils ne conservent pas le
quasi-moment, donc ils dégradent le flux : \textbf{ils sont la source de la résistance thermique
intrinsèque d'un cristal parfait}.

Ils exigent des phonons de grand vecteur d'onde, donc de haute énergie, et sont par conséquent gelés
à basse température. C'est la raison pour laquelle la conductivité thermique d'un cristal pur croît
fortement lorsque la température diminue, avant d'être limitée par d'autres mécanismes.

\subsection{Autres mécanismes résistifs}

La diffusion par les impuretés et les isotopes est regroupée avec les processus umklapp dans
un temps résistif volumique $\tau_R$ du modèle gris. Les frontières interviennent par
leurs conditions de bord et par un temps géométrique indicatif $\tau_B$ ; elles ne sont
pas ajoutées une seconde fois dans ce $\tau_R$. Leur résistance dépend notamment de la spécularité.

Dans un film mince, la diffusion aux frontières prend une importance particulière : dès que
l'épaisseur devient comparable aux longueurs de transport pertinentes, elle peut réduire la
conductivité apparente ; son importance exige le spectre et les propriétés des frontières.

\subsection{Critère de Guyer}

Le comportement du gaz de phonons est fixé par la hiérarchie des fréquences de collision. En notant
$\tau_B$ le temps associé à la diffusion aux frontières, la fenêtre dite hydrodynamique est définie
par
\[ \frac{1}{\tau_N} \;\gg\; \frac{1}{\tau_B} \;\gg\; \frac{1}{\tau_R} . \]

Les processus normaux doivent être très fréquents, pour que le gaz reste en équilibre interne, et
les processus résistifs très rares, pour que le flux ne soit pas dégradé. Entre les deux, la
géométrie fixe l'échelle.

Hors de cette fenêtre, deux autres régimes apparaissent : le régime diffusif lorsque les processus
résistifs dominent, et le régime balistique lorsque les frontières dominent tout le reste.

\section{L'approximation de Callaway}

\subsection{Principe}

Le terme de collision exact est intraitable. L'approximation de Callaway le remplace par une
relaxation vers deux états d'équilibre distincts, chacun avec son temps propre :
\[ \boxed{\;\left(\frac{\partial f}{\partial t}\right)_{\!\mathrm{coll}}
= -\,\frac{f - f^{\,\mathrm{d}}}{\tau_N} \;-\; \frac{f - f^{\,0}}{\tau_R}\;} \]

Le premier terme décrit les processus normaux. Ils ramènent la distribution vers la forme déplacée
$f^{\,\mathrm{d}}$, qui porte une vitesse d'ensemble $\mathbf{u}$ non nulle : le gaz s'équilibre
sans cesser de s'écouler.

Le second décrit les processus résistifs. Ils ramènent la distribution vers l'équilibre immobile
$f^{\,0}$, de vitesse d'ensemble nulle : le gaz est arrêté.

\subsection{Ce que l'approximation suppose}

Elle suppose que le retour à l'équilibre est exponentiel, avec un temps caractéristique unique par
famille de processus. Cette hypothèse efface la dépendance en fréquence des temps de collision, qui
est en réalité marquée.

Les calculs de premiers principes sur l'AlN wurtzite donnent, à basse fréquence,
\[ \frac{1}{\tau_U} \propto \omega^3 \quad\text{(TA et LA)} , \qquad
\frac{1}{\tau_N} \propto \omega \quad\text{(TA)} , \qquad
\frac{1}{\tau_N} \propto \omega^2 \quad\text{(LA)} . \]

Le rapport $\tau_R/\tau_N$ varie donc comme $\omega^{-2}$ pour les modes transverses : il n'est en
aucune façon constant sur le spectre.

Les temps constants employés ensuite définissent un modèle gris. Leur interprétation comme
paramètres effectifs exige une pondération dérivée de la fermeture, et non une moyenne choisie
arbitrairement. La note~14 précise cette restriction ; aucune moyenne spectrale de l'AlN n'est fournie.

\subsection{Ce qu'elle préserve}

Elle respecte les deux lois de conservation qui importent. L'énergie est conservée par les deux
familles de processus, ce qui impose que les deux termes de collision s'annulent une fois
multipliés par $\hbar\omega$ et intégrés. Le quasi-moment est conservé par les seuls processus
normaux, ce qui impose que le premier terme s'annule une fois multiplié par $\hbar\mathbf{k}$ et
intégré.

Ces deux propriétés sont exactement ce qui produit la structure des équations de moments.

\section{La hiérarchie de moments}

\subsection{Moment d'ordre zéro}

On multiplie l'équation de transport par $\hbar\omega_\mathbf{k}$ et on intègre sur les vecteurs
d'onde. Le premier terme donne $\partial u/\partial t$, le second $\nabla \cdot \mathbf{q}$, et le
terme de collision s'annule par conservation de l'énergie. Il reste
\[ \boxed{\;\frac{\partial u}{\partial t} + \nabla \cdot \mathbf{q} = 0\;} \]

C'est le bilan d'énergie, obtenu sans aucune hypothèse supplémentaire. En écrivant
$\partial u/\partial t = C\,\partial T/\partial t$, on retrouve l'équation employée dans tout le
travail de modélisation.

Cette équation est exacte mais incomplète : elle relie deux inconnues, $u$ et $\mathbf{q}$, et n'en
détermine aucune.

\subsection{Moment d'ordre un}

On multiplie cette fois par $\hbar\omega_\mathbf{k}\,\mathbf{v}_g$ et on intègre. Le premier terme
donne $\partial \mathbf{q}/\partial t$. Le second fait apparaître un moment d'ordre supérieur,
\[ \mathbf{Q} = \int \frac{d^3k}{(2\pi)^3}\; \hbar\omega\; \mathbf{v}_g \otimes \mathbf{v}_g\; f , \]
tenseur d'ordre deux qui décrit le transport du flux lui-même.

Pour le terme de collision, la proportionnalité $\mathbf{q} = v^2 \mathbf{P}$ établie dans les notes
précédentes est décisive : les processus normaux conservant $\mathbf{P}$, ils conservent
$\mathbf{q}$ et ne contribuent pas. Seuls les processus résistifs subsistent, et
\[ \boxed{\;\frac{\partial \mathbf{q}}{\partial t} + \nabla \cdot \mathbf{Q}
= -\,\frac{\mathbf{q}}{\tau_R}\;} \]

\subsection{Le problème de fermeture}

Le schéma se répète indéfiniment. L'équation du moment d'ordre zéro appelle le moment d'ordre un ;
celle du moment d'ordre un appelle le moment d'ordre deux ; et ainsi de suite.

\begin{center}
\begin{tabular}{@{}lll@{}}
\toprule
Ordre & Équation obtenue & Inconnue introduite \\
\midrule
0 & bilan d'énergie & le flux $\mathbf{q}$ \\
1 & évolution du flux & le tenseur $\mathbf{Q}$ \\
2 & évolution de $\mathbf{Q}$ & le moment d'ordre trois \\
$\vdots$ & $\vdots$ & $\vdots$ \\
\bottomrule
\end{tabular}
\end{center}

Aucun niveau ne se referme de lui-même. Obtenir un système fini exige d'exprimer un moment en
fonction des précédents, par une hypothèse qui ne découle pas de l'équation de transport. C'est la
\emph{fermeture}, et c'est là que réside tout le contenu physique de l'approximation.

\section{Fermeture et convention constitutive}
La dérivation complète est désormais donnée dans la note~14. Dans le modèle
gris linéarisé, $Q_{ij}=Cv^2\theta\,\delta_{ij}/3+\Pi_{ij}$ et
$\Pi_{ii}=0$. Avec
$\tau_c=(\tau_N^{-1}+\tau_R^{-1})^{-1}$, l'élimination du moment rapide donne
\[
 \Pi_{ij}=-\frac{v^2\tau_c}{5}
 \left(\partial_iq_j+\partial_jq_i-\frac23\delta_{ij}\nabla\cdot\mathbf q\right).
\]
La divergence du tenseur conduit à
\[
 \tau_R\partial_t\mathbf q+\mathbf q=-\lambda\nabla T+
 \ell_c^2\left[\nabla^2\mathbf q+\frac13\nabla(\nabla\cdot\mathbf q)\right],
 \quad
 \lambda=\frac13Cv^2\tau_R,\quad \ell_c^2=\frac15v^2\tau_R\tau_c .
\]
La limite collective donne $\tau_c\simeq\tau_N$.
Elle exige $\tau_N/\tau_R\ll1$, $|p|\tau_N\ll1$ et $v\tau_N/L\ll1$.
La projection sans trace préserve la définition énergétique de la température.

La forme GK historique remplace $1/3$ par $2$. Elle est utilisée dans les
notes antérieures comme modèle constitutif, mais n'est pas le résultat de
cette fermeture conservatrice. La différence est discutée dans Sendra et al.,
Phys. Rev. B \textbf{106}, 155301 (2022), équation (21).
Les deux conventions ont les mêmes limites Fourier et Cattaneo.
En 1D on pose $L^2=(1+\alpha)\ell^2$ et $\tau_\ell=L^2/a$ :
$\alpha=2$ pour GK historique, $\alpha=1/3$ pour la fermeture conservatrice.
Le code utilise directement $\tau_\ell$, sans imposer une interprétation de $\ell$.

\section{Les limites}

\subsection{Limite de Cattaneo}

Lorsque la longueur non locale devient négligeable devant l'échelle de variation du flux, les termes
en $\ell^2$ disparaissent et il reste
\[ \tau_R\,\frac{\partial \mathbf{q}}{\partial t} + \mathbf{q} = -\,\lambda\,\nabla T . \]

C'est l'équation de Cattaneo, à un seul temps de relaxation. Pour $\tau_R>0$,
elle est strictement hyperbolique et sa vitesse caractéristique est finie.
Le terme de diffusion spatiale du flux a disparu.

\subsection{Limite de Fourier}

Lorsque le temps de relaxation devient en outre négligeable devant l'échelle de temps de
l'expérience, le terme instationnaire disparaît et il reste
\[ \mathbf{q} = -\,\lambda\,\nabla T . \]

C'est la loi de Fourier, dont la combinaison avec le bilan d'énergie donne l'équation de la chaleur
parabolique, de vitesse de propagation infinie.

\subsection{Récapitulation}

\begin{center}
\begin{tabularx}{\textwidth}{@{}lXl@{}}
\toprule
\textbf{Modèle} & \textbf{Condition} & \textbf{Paramètres} \\
\midrule
Guyer--Krumhansl & fenêtre hydrodynamique de Guyer & $\lambda$, $\tau_R$, $\ell$ \\
\addlinespace
Cattaneo & $\ell$ négligeable devant l'échelle spatiale & $\lambda$, $\tau_R$ \\
\addlinespace
Fourier & $\tau_R$ négligeable devant l'échelle temporelle & $\lambda$ \\
\bottomrule
\end{tabularx}
\end{center}

\section{Les deux paramètres et leurs dimensions}

\subsection{Un temps et une longueur}

L'équation de Guyer--Krumhansl introduit deux grandeurs au-delà de la conductivité.

Le temps résistif $\tau_R$, en secondes, est le temps de relaxation qui figure dans la loi de
Cattaneo. Sa valeur dépend du matériau et de la moyenne cinétique adoptée ;
les temps des modes individuels ne sont pas un temps effectif unique.

La longueur non locale $\ell = v\sqrt{\tau_N \tau_R / 5}$, en mètres, n'a pas d'équivalent dans le
modèle de Cattaneo. Sa valeur doit être calculée avec des temps compatibles
avec cette fermeture, et ne doit pas être assimilée à une longueur grise issue de la conductivité.

\subsection{Conséquence pour l'invariance d'échelle}
Dans le modèle homogène 1D du projet, la réponse dépend de
$b,\xi_1,\tau_R$ et $\tau_\ell=L^2/a$, avec $L^2=3\ell^2$ dans la convention historique.
La transformation $d\mapsto\mu d$, $\lambda\mapsto\mu\lambda$,
$C\mapsto C/\mu$, $\ell\mapsto\mu\ell$ conserve ces quatre quantités.
La longueur non locale ne lève donc pas la dégénérescence lorsque ces paramètres sont libres.

Si $v$ est connu indépendamment et si $a=v^2\tau_R/3$ est imposé,
cette information supplémentaire contraint $a$ ; l'invariance précédente
ne respecte alors plus toutes les données imposées.
L'extension aux profils gradués non-Fourier reste hors du code actuel.

\section{Limites du cadre}

\subsection{Erreur propre de l'approximation de Callaway}

Une comparaison entre le modèle de Callaway et la résolution exacte de l'équation de Boltzmann, menée
par méthode itérative sur des taux de diffusion calculés de premiers principes, conclut que ni le
modèle original ni la version modifiée d'Allen ne garantissent une amélioration sur l'approximation
du temps de relaxation.

Pour l'AlN wurtzite, l'écart du modèle de Callaway à la solution exacte vaut $+1\,\%$ dans le plan et
$-12\,\%$ hors plan.

\textbf{La direction hors plan est celle du présent travail.}
L'écart cité concerne une conductivité stationnaire calculée, pour le modèle
et le matériau de Ma, Li et Luo (2014). Il n'est pas une incertitude universelle
de dix pour cent sur une réponse dynamique.

\subsection{Portée}
Un temps de relaxation extrait ne reçoit pas automatiquement la même erreur
relative que la conductivité. Les erreurs de fermeture doivent être propagées
à l'observable et à l'estimateur. Un désaccord expérimental de moins de dix pour cent
ne peut pas être déclaré acceptable sans cette analyse et les erreurs de mesure.

\section{Références}

{\small
\begin{enumerate}[leftmargin=1.6em,itemsep=0pt]
\item J. Callaway, \emph{Model for lattice thermal conductivity at low temperatures}, Phys. Rev.
\textbf{113}, 1046 (1959).
\item R. A. Guyer, J. A. Krumhansl, \emph{Solution of the linearized phonon Boltzmann equation},
Phys. Rev. \textbf{148}, 766 (1966).
\item R. A. Guyer, J. A. Krumhansl, \emph{Thermal conductivity, second sound, and phonon
hydrodynamic phenomena in nonmetallic crystals}, Phys. Rev. \textbf{148}, 778 (1966).
\item N. W. Ashcroft, N. D. Mermin, \emph{Solid State Physics}, Holt (1976).
\item C. Kittel, \emph{Introduction to Solid State Physics}, Wiley.
\item G. P. Srivastava, \emph{The Physics of Phonons}, Adam Hilger (1990).
\item Examen du modèle de Callaway pour Si, diamant et AlN wurtzite, Phys. Rev. B \textbf{90},
035203. Source des lois d'échelle en fréquence et de l'écart hors plan.
\item Mesures de conductivité sur AlN monocristallin de 130 à 480 K, Phys. Rev. Materials
\textbf{4}, 044602.
\end{enumerate}}

\vspace{4pt}
\noindent\small
La fermeture grise et ses contrôles sont détaillés dans la note~14 ; la comparaison des sources et ses limites d'accès figurent dans la note~15. Notes associées : \texttt{08\_\allowbreak Le\_\allowbreak gaz\_\allowbreak de\_\allowbreak phonons},
\texttt{07\_\allowbreak Note\_\allowbreak partie\_\allowbreak IV}.
\end{document}